\section{Projectives over finite distributive support extensions}
\label{sec:finite-lattice-projectives}
We prove Theorem~\ref{thm:finite-lattice-projectives}. Throughout this
section $C$ is a finite poset, $L=\Id(C)$, and
$T=G_L(R)$ for a unital ring $R$, possibly noncommutative. The Boolean
projectivity input is the finite case of Horn and Kimura's theorem,
proved in Lemma~\ref{lem:lattice}; see also Macpherson
\cite[Theorem~3.9]{Macpherson}. The group-diagram argument below is
ordinary projective-module theory on a finite poset.

\begin{lemma}[module data]\label{lem:lattice-module-data}
A $T$-semimodule is equivalently an $L$-semimodule $N$ together with a
covariant diagram $G:N\to R\text{-}\Mod$. Morphisms consist of an
$L$-linear support map and compatible linear maps on groups. The fiber
over the least element of $N$ need not be zero.
\end{lemma}
\begin{proof}
The action of $(-1_R,C)$ supplies the additive inverses in the sense
of inverse monoids. Their group identities form $N$. Every
$(0_R,l)$ acts on identities as the corresponding lattice scalar;
$(0_R,C)$ acts as their identity. The old ring elements act within
each group, making it an $R$-module. Adding a larger group identity
gives the transition maps, as in Section~\ref{sec:diagrams}.

Conversely, transport addition to the join group, let an old ring
scalar act linearly in the same fiber, and let a proper lattice scalar
give the zero in its scalar-translated fiber. These formulas satisfy
the module identities. In particular, if two proper lattice elements
join to $C$, their sum in $T$ is $(0_R,C)$, the old zero; both sides of
the distributive identity are the group zero in the join fiber.
The zero scalar gives the zero in the least fiber. The constructions
on maps are inverse because a homomorphism preserves group identities
and additive inverses.
\end{proof}

\begin{lemma}[the support classification]\label{lem:lattice-support-projectives}
The finite projective $L$-semimodules are exactly the colored order-ideal
lattices in the first part of
Theorem~\ref{thm:finite-lattice-projectives}.
\end{lemma}
\begin{proof}
Write $N$ as a retract of $L^n$, with section $i$ and retraction $r$.
Since $L$ is a finite projective Boolean semimodule, so is $N$;
Lemma~\ref{lem:lattice} gives $N=\Id(P)$ for a finite poset $P$.
Let $d_l$ denote multiplication by $l\in L$. If $I\le J$ and
$d_lJ=J$, then $i(I)\le i(J)$ has every coordinate contained in $l$,
so $d_lI=I$. Thus the fixed elements are the ideals contained in a
fixed ideal $U_l\subseteq P$. The map $d_l$ sends $I$ to its greatest
$l$-fixed subideal, hence $d_lI=I\cap U_l$. The action laws give
\[
 U_{l\cap m}=U_l\cap U_m,\quad U_{l\cup m}=U_l\cup U_m,
 \quad U_\varnothing=\varnothing,\quad U_C=P.
\]
For each $p$, membership in $U_l$ is a bounded lattice character of
$\Id(C)$, hence evaluation at a unique $c(p)\in C$. The ideals $U_l$
make $c$ order preserving, and give the asserted scalar action.

Decompose $i(\mathord\downarrow p)$ into coordinate join-irreducibles
$(\mathord\downarrow a)e_j$ of $L^n$. Some term retracts to
$\mathord\downarrow p$. Its color is $a=c(p)$: the section of
$\mathord\downarrow p$ is fixed by $\mathord\downarrow c(p)$,
giving $a\le c(p)$, while its retraction is fixed by
$\mathord\downarrow a$, giving the reverse inequality. Fix that
coordinate $j$. For $a\le c(p)$ put
\[
 X=r((\mathord\downarrow a)e_j)
   =d_{\mathord\downarrow a}(\mathord\downarrow p).
\]
The $j$th coordinate of $i(X)$ is exactly $\mathord\downarrow a$.
If $X=Y\vee Z$, the join-primality of $a$ in $L$ puts $a$ in one
of the corresponding section coordinates, say that of $Y$.
Then $(\mathord\downarrow a)e_j\le i(Y)$ and $X\le Y$.
Therefore $X$ is a nonzero join-irreducible of $N$. Its color is exactly
$a$: it is fixed by $\mathord\downarrow a$, and every scalar fixing
it must fix its $j$th section coordinate. Thus $X=\mathord\downarrow p_a$
with $c(p_a)=a$, proving the required greatest-element condition.

Conversely, for such a colored poset define
\begin{equation}\label{eq:lattice-support-section}
 s(I)_p=\{a\le c(p):p_a\in I\},\qquad
 r(l e_p)=l\cdot\mathord\downarrow p.
\end{equation}
The greatest-element property gives $(p_b)_a=p_a$ for
$a\le b\le c(p)$, and $p_a\le q_a$ when $p\le q$.
These identities make $s(I)_p$ an ideal in $C$. Membership of $p_a$
proves that $s$ preserves zero and joins, and its exact color proves
$s(lI)=l s(I)$. Every term of $r(s(I))$ lies in $I$, while $p\in I$
implies $c(p)\in s(I)_p$ and hence $p\in r(s(I))$. Thus $rs=1_N$.
\end{proof}

\begin{lemma}[projective diagrams on a finite poset]
\label{lem:projective-poset-diagrams}
Let $X$ be a finite poset. Every finitely generated projective object
$G$ of $\operatorname{Fun}(X,R\text{-}\Mod)$ is isomorphic to
\[
 G_x=\bigoplus_{a\le x}Q_a
\]
with inclusion maps and finitely generated projective $R$-modules $Q_a$.
Their isomorphism classes are recovered by
\begin{equation}\label{eq:diagram-birth-quotient}
 Q_a\simeq G_a\big/\sum_{b<a}\im(G_b\longrightarrow G_a).
\end{equation}
\end{lemma}
\begin{proof}
Let $E_a(Q)$ be the diagram equal to $Q$ at $x\ge a$ and zero elsewhere,
with identity maps wherever both values are $Q$. The diagrams $E_a(R)$
are projective generators: maps from them are elements of the group
at $a$. Therefore $G$ is a retract of a finite sum of these diagrams.
The quotient functor on the right of
\eqref{eq:diagram-birth-quotient} sends $E_b(R)$ to $R$ when $a=b$
and to zero otherwise. It preserves direct sums and retracts, so its
value $Q_a$ on $G$ is finitely generated projective over $R$.
Choose a linear section $Q_a\to G_a$ of each quotient. Transporting
these sections gives a natural map $F=\bigoplus_aE_a(Q_a)\to G$.
It is pointwise surjective by induction on $X$, and splits because
$G$ is projective. Applying each quotient functor to the splitting
shows that its kernel has every such quotient zero. A second induction
on $X$ makes the kernel zero. This proves both assertions.
\end{proof}

\begin{proof}[Proof of Theorem~\ref{thm:finite-lattice-projectives}]
Let $M$ be a finitely generated projective $T$-semimodule, with support
$N$ and group diagram $G$. The support is projective over $L$.
Indeed, lift the support map $M\to N$ across a surjection of $L$-modules,
regarding both target and domain as zero-group $T$-modules, and restrict
the lift to group identities. It is finite because it is finitely
generated over the finite semiring $L$. Lemma~\ref{lem:lattice-support-projectives}
gives $N=\Id(P)$ with the stated coloring.

The diagram $G$ is projective in $\operatorname{Fun}(N,R\text{-}\Mod)$.
Any pointwise surjection $H\to G$ defines, by
Lemma~\ref{lem:lattice-module-data}, a surjective $T$-module map with
support identity. A section supplied by projectivity of $M$ also has
support identity, so is a diagram section. Moreover a finite generating
set of $M$ generates $G$ as a diagram: all group contributions are
linear combinations of transported generators, since proper lattice
scalars contribute only group zeros. Lemma~\ref{lem:projective-poset-diagrams}
therefore gives labels $Q_A$ for all $A\in N$.

We show that a nonzero label requires admissibility. Choose a free
retraction $M\longrightarrow T^n\longrightarrow M$, with support
maps $i$ and $r$, and put $A_j=r(Ce_j)$. The group fiber of $T^n$
at $i(A)$ has free coordinates exactly at indices with $i(A)_j=C$.
Each such coordinate satisfies $A_j\subseteq A$, and its retraction
comes from the group $G_{A_j}$. If every active coordinate had
$A_j\subsetneq A$, the identity composite on $G_A$ would induce zero
on the quotient \eqref{eq:diagram-birth-quotient}, forcing $Q_A=0$.
Consequently $Q_A\ne0$ gives a coordinate $j$ with
$i(A)_j=C$ and $A_j=A$.

Set $\psi(I)=i(I)_j$ and $\theta(l)=lA$. They satisfy
\[
 \psi\theta=1_L,\qquad \theta\psi(I)\subseteq I.
\]
For $a\in C$, the element $X=\theta(\mathord\downarrow a)$ has
$\psi(X)=\mathord\downarrow a$. If $X=Y\vee Z$, membership of $a$
in one of $\psi(Y),\psi(Z)$ forces $X\le Y$ or $X\le Z$.
Thus $X$ is nonzero join-irreducible. It is fixed by
$\mathord\downarrow a$, while any scalar fixing $X$ must fix
$\psi(X)$; its color is therefore exactly $a$.
This is \eqref{eq:admissible-ideal}.

Conversely, start from the data in the statement. For each admissible
$A$, let $p_a$ denote its color-cut maximum and define
\begin{equation}\label{eq:admissible-functional}
 \rho_A(I)=\{a\in C:p_a\in I\}.
\end{equation}
It is a zero-preserving $L$-linear join map. It has
$\rho_A(I)=C$ exactly when $A\subseteq I$, and
$\rho_A(I)A\subseteq I$. Start with the support retraction
\eqref{eq:lattice-support-section}, giving every one of those
coordinates zero amplitude. For each nonzero $Q_A$, choose split
linear maps $Q_A\to R^{m_A}\to Q_A$ and append $m_A$ coordinates.
Their support section is $\rho_A$, and each corresponding retraction
column sends full support to $A$. On groups use the chosen module
section and retraction for the $Q_A$ summand. Outside $I\supseteq A$
the corresponding coordinate has proper support and zero amplitude,
so these formulas define $T$-linear maps.

The support composite is the identity: the original coordinates
already recover $I$, and every appended contribution lies in
$\rho_A(I)A\subseteq I$. The group composite is the identity by
the chosen module splittings. Thus the constructed module is a
retract of a finite free one. Finally $N$ recovers $P$ and its coloring,
and \eqref{eq:diagram-birth-quotient} recovers every label. This proves
the isomorphism classification. Diagram isomorphisms may include
upper-triangular mixing of summands; only their isomorphism classes
are being classified.
\end{proof}

\begin{corollary}[Grothendieck groups and strong projectives]
\label{cor:lattice-projective-Kzero}
Assume $C\ne\varnothing$. For data in
Theorem~\ref{thm:finite-lattice-projectives}, join two vertices of $P$
when they are comparable, and join every pair of vertices belonging
to an admissible $A$ with $Q_A\ne0$. Let $\mathcal D_{C,R}$ be the set
of isomorphism classes of nonempty data connected in this graph.
Then
\[
 V(G_{\Id(C)}(R))\simeq\N^{(\mathcal D_{C,R})},\qquad
 K_0(G_{\Id(C)}(R))\simeq\Z^{(\mathcal D_{C,R})}.
\]
If $R\ne0$, every finitely generated strong projective is free, and
$SV(G_{\Id(C)}(R))\simeq\N$, $SK_0(G_{\Id(C)}(R))\simeq\Z$.
If $R=0$, this monoid and group are instead
$\N^{\pi_0(C)}$ and $\Z^{\pi_0(C)}$, respectively, where $\pi_0(C)$
is the set of comparability components of $C$. When $C=\varnothing$
all assertions reduce to the corresponding ordinary ring invariants.
\end{corollary}
\begin{proof}
A direct sum takes the disjoint union of the colored posets, retains
the labels on ideals wholly contained in a summand, and gives zero
labels to every mixed admissible ideal. Since $C$ is nonempty,
every admissible ideal is nonempty. A decomposition therefore amounts
exactly to a partition into comparability components which does not
separate any nonzero label. The displayed graph gives the unique
finest partition. Restricting data to any component retains the
color-cut condition, because its required greatest elements are
comparable to their original vertices; all its nonzero labels remain
admissible. This proves the free-monoid assertion and its group
completion.

The free rank-one semimodule has $P=C$, identity coloring, and the
single label $R$ on $A=C$. If $R\ne0$, this label connects all of $C$.
A direct summand of its $n$-fold direct sum is therefore a union of
whole copies of these data, hence free. If $R=0$, the labels vanish
and the indecomposable pieces of the free rank-one object are the
comparability components of $C$. They remain distinguishable even
if their uncolored posets are isomorphic, because colors are fixed.
An arbitrary finite multiplicity for each such component is a direct
summand of a sufficiently large free module, giving the stated
monoid and group. The empty-$C$ case is $T=R$.
\end{proof}
